% prop-norm-shape -- CERW.Frozen.norm_shape_rates (draft). Source: limit-shapes.tex:746-762, label prop:norm-shape
\paragraph{Paper (verbatim).}
\begin{proposition}\label{prop:norm-shape}
Let $\gauge$ be a norm, let $\xi(x)\in\partial \gauge(x)$ be subgradients chosen at the sites $x\in\Z^d\setminus\{0\}$, and let $\drift>0$ satisfy~\eqref{eq:ellipticity}. For every $p>0$ there exists $C(d,\drift,\gauge,p)<\infty$, not depending on the choice of the subgradients, such that, for all integers~$n\geq2$, with probability at least $1-Cn^{-p}$ the following hold simultaneously.
\begin{enumerate}[label=\textup{(\roman*)}]
\item \underline{\emph{Radii}}:
\begin{equation}\label{eq:norm-shape}
\Bigl|\inf_{y\notin D_n}\gauge(y)-r_n\Bigr|\leq C\begin{cases}r_n^{\nf34}(\log n)^{\nf14},&d=2,\\(r_n\log n)^{\nf12},&d\geq3\end{cases}
\end{equation}
and
\begin{equation}\label{eq:norm-outer}
\max_{0\leq j\leq n}\gauge(X_j)-r_n\leq C\begin{cases}r_n^{\nf56}(\log n)^{\nf76},&d=2,\\r_n^{\nf{d}{(d+1)}}(\log n)^{\nf{(d+2)}{(d+1)}},&d\geq3\end{cases}\, .
\end{equation}
\item \underline{\emph{Local times}}:
\begin{equation}\label{eq:norm-local}
\max_{x\in\Z^d}|\ell_n(x)-2d\drift(r_n-\gauge(x))_+|\leq C\begin{cases}r_n^{\nf56}(\log n)^{\nf16},&d=2,\\r_n^{\nf{d}{(d+1)}}(\log n)^{\nf{1}{(d+1)}},&d\geq3\end{cases}\, .
\end{equation}
\end{enumerate}
\end{proposition}
Definitions the statement relies on (verbatim): $r_n$ (lines 323--325), $D_n,\Rin,\Rout$ (lines 122--127).
\begin{equation}\label{eq:radius-norm}
r_n\coloneqq\left(\frac{(d+1)n}{2d\drift|B_\gauge|}\right)^{\nf{1}{(d+1)}}\, ,
\end{equation}
Replace each site~$x\in A_n$ by the unit cell~$C_x\coloneqq x+[-\nf12,\nf12)^d$. Let $D_n\coloneqq\bigcup_{x\in A_n}C_x$. The \defn{inner radius} and the \defn{outer radius} are
\begin{equation}\label{eq:radii}
\Rin(n)\coloneqq\inf_{y\notin D_n}|y|
\qquad\text{and}\qquad
\Rout(n)\coloneqq\max_{0\leq j\leq n}|X_j|\, .
\end{equation}
\paragraph{Transcription.}
\begin{itemize}
\item \emph{Setting and constants.} As in the coarse-bounds twin. One constant $C(d,\drift,\gauge,p)$, bound after $p$ and before
 $\xi$ and the walk (``not depending on the choice of the subgradients''), serves in the three bounds and in the failure
 probability.
\item \emph{Case split.} $d=2$ versus $d\ge3$ is \texttt{if d = 2 then \dots else \dots} ($d\ge2$ is \texttt{hd}). Exponents, with
 \texttt{R}-valued real powers: $d=2$: $r^{3/4}(\log n)^{1/4}$, $r^{5/6}(\log n)^{7/6}$, $r^{5/6}(\log n)^{1/6}$. $d\ge3$:
 $(r\log n)^{1/2}$, $r^{d/(d+1)}(\log n)^{(d+2)/(d+1)}$, $r^{d/(d+1)}(\log n)^{1/(d+1)}$. They match eq:norm-shape, eq:norm-outer,
 eq:norm-local.
\item \emph{(i) Inner radius.} $\inf_{y\notin D_n}\gauge(y)=$ \texttt{normInnerRadius Psi (X . w) n} $=\texttt{sInf}\,(\gauge''(D_n)^{\complement})$
 with $D_n=$ \texttt{cellSet} the union of the half-open unit cells of $A_n$; the image is nonempty ($D_n$ bounded) and bounded
 below by $0$, so \texttt{sInf} is the infimum. The claim is $|\text{that}-r_n|\le\dots$.
\item \emph{(i) Outer radius.} $\max_{0\le j\le n}\gauge(X_j)=$ \texttt{normMaxRadius Psi (X . w) n} ($\sup'$ over $j\in\{0,\dots,n\}$ of
 $\gauge(\texttt{toSpace}(X_j))$); the claim is the one-sided bound $\text{that}-r_n\le\dots$.
\item \emph{(ii) Local times.} $\forall x:\texttt{Site d}$, $|\ell_n(x)-2d\drift\max(r_n-\gauge(x),0)|\le\dots$, where $(r_n-\gauge(x))_+$
 is \texttt{max (r n - Psi (toSpace x)) 0} and $\ell_n=$ \texttt{localTime}. The maximum over $x\in\Z^d$ is the universal quantifier
 inside the event (for $|x|$ large both terms vanish).
\item \emph{Simultaneity and probability.} The three parts are one conjunction inside the event; the bound is
 $\mu\{\omega:\neg(\cdots)\}\le\texttt{ENNReal.ofReal}(C\,n^{-p})$ for every $n\ge2$.
\end{itemize}
